% Method.
\section{Method: GA-DCHG}
\label{sec:method}

\subsection{Phase-space hypergraph baseline}
\label{sec:baseline}

From a vibration signal $x(t)$ we reconstruct the phase space by time-delay
embedding \cite{takens1981,fraser1986,kennel1992}, yielding a trajectory of
state points $\bm{x}_t\in\mathbb{R}^m$; phase-space reconstruction and its
geometry have been used for rotor fault diagnosis in complex networks
\cite{zhao2025icmds}. The trajectory is partitioned into
segments (windows); within a segment the \emph{dynamical distance}
\begin{equation}
  d^{dyn}_{ij} \;=\; \lVert \bm{x}_i - \bm{x}_j \rVert \cdot \exp\!\bigl(
  -\tfrac{|t_i - t_j|}{\sigma}\bigr)
\label{eq:dyn}
\end{equation}
combines geometric separation with a time-decay that prefers temporally local
neighbours, and a $k$-nearest-neighbour rule on $d^{dyn}$ defines hyperedges
aggregated by a hypergraph layer \cite{feng2019}. The baseline gates its
attention with three \emph{edge-geometry descriptors} of each hyperedge $e$
\cite{grassberger1983}: local curvature $\kappa_e$, expansion $\chi_e$, and
filling $\rho_e$, collected as $\bm{\phi}_e=(\kappa_e,\chi_e,\rho_e)$. These
descriptors are \emph{passive}: they gate attention but do not shape how
hyperedges are formed.
As shown in Sec.~\ref{sec:results}, the phase-space geometry itself is markedly
more invariant to a load shift than the raw waveform, so the structural model
already transfers far better than a raw 1D-CNN under load. GA-DCHG exploits this
by making the geometry \emph{active}: the descriptors become anchors that decide
which hyperedges are generated, how they evolve over time, and how source and
target are aligned.

\subsection{C1: data-driven anchor screening}
\label{sec:c1}

Which descriptor is stable across a domain shift is answered empirically, not by
assumption. For a protocol with source condition $\mathcal{S}$ and target
condition $\mathcal{T}$, we pool the per-edge descriptor values in each domain
and measure, for each $d\in\{\kappa,\chi,\rho\}$, the normalized
$1$-Wasserstein distance $W_1(d)=\mathrm{Wass}_1(p_d^{\mathcal{S}},
p_d^{\mathcal{T}})/\mathrm{std}(p_d^{\mathcal{S}}\cup p_d^{\mathcal{T}})$
between the source and target pools, using only unlabeled target windows. The
\emph{geometric anchor} is the descriptor with the smallest distance,
$a^\star=\arg\min_d W_1(d)$ (Table~\ref{tab:screen}); the screening gives a
per-protocol anchor, not a fixed descriptor (Sec.~\ref{sec:r1}).

\subsection{C2: anchor-modulated dynamic hyperedge generation}
\label{sec:c2}

For a window we compute the anchor descriptor $h_e$ of each hyperedge
$e$ and its mean $\alpha_s$ over the hyperedges of segment $s$. Because the fault
trajectory is smooth, the anchor state is propagated along the segment index by
an exponential moving average,
\begin{equation}
  \bar\alpha_s \;=\; \tau\,\bar\alpha_{s-1} + (1-\tau)\,\alpha_s,
\label{eq:ema}
\end{equation}
giving each hyperedge a short-term memory of geometric drift. The modulated
distance used for hyperedge formation is
\begin{equation}
  \tilde d_{ij} \;=\; d^{dyn}_{ij}\,\Bigl(1 + \eta\,\bigl|\tfrac{1}{2}
  (h_i + h_j) - \bar\alpha_s\bigr|\Bigr),
\label{eq:mod}
\end{equation}
which strengthens connections between pairs whose anchor geometry is close to
the running trajectory state and weakens those that deviate. Temporal
propagation is applied \emph{only} to the anchor descriptor; non-anchor
descriptors remain window-local so that domain-specific geometry cannot pollute
the cross-domain state.

\subsection{C3: geometry-space alignment}
\label{sec:c3}

Instead of adversarial embedding alignment, we align the \emph{effective anchor
geometry}: the attention-weighted anchor statistic $g_{\bm w}$ of a window
(its learned hyperedge attention weights $\beta_{\bm w e}$, defined as in
\cite{vaswani2017}) is a scalar, and the alignment loss is the closed-form
$1$-Wasserstein distance between the source and target-adapt sets of these
scalars,
\begin{equation}
  \mathcal{L}_{\mathrm{align}} \;=\; \mathrm{Wass}_1\!\bigl(\{g_{\bm w}\}_{\bm w\in
  \mathcal{S}}, \{g_{\bm w}\}_{\bm w \in \mathcal{T}_{\mathrm{adapt}}}\bigr),
\label{eq:align}
\end{equation}
added to the source cross-entropy with weight $\lambda$:
$\mathcal{L}=\mathcal{L}_{\mathrm{CE}}+\lambda\,\mathcal{L}_{\mathrm{align}}$.
Because the anchor is a scalar, the $1$-Wasserstein term needs only a sort;
it equates how the fault trajectory fills the phase space in the two domains and
is computationally free.

\subsection{C4: strict protocol and intervention audit}
\label{sec:c4}

Source and target never share a recording; the target is split into an unlabeled
\emph{adaptation} set (65\%) and a held-out \emph{test} set (35\%), also
recording-disjoint. We train on labeled source windows plus the target-adapt set
(for alignment) and report target-test accuracy. To evidence that the decision
\emph{uses} the geometry, we perturb descriptor $j$ at test time,
$\phi^{(j)}_e \leftarrow \phi^{(j)}_e + \delta\,\mathrm{std}_j\,z_e$ with
$z_e\sim\mathcal{N}(0,1)$, and measure the \emph{flip rate}, defined as the
fraction of test windows whose prediction changes. A decision that uses the geometry is
measurably disrupted by the perturbation, whereas an ignored descriptor leaves
decisions unchanged.
